\documentclass{report}
\usepackage{amsfonts}
\newcommand\R{{\mathbb R}}
\newcommand\Q{{\mathbb Q}}
\newcommand\N{{\mathbb N}}
\pagestyle{empty}
\begin{document}
\large
\centerline{\Huge Analisi Matematica}
\renewcommand{\thefootnote}{ } 
\footnote{\sl Avete 3 ore di tempo. Ogni esercizio
vale sei punti. La sufficienza si ottiene con un punteggio $\geq$
18. Solo le risposte chiaramente giustificate saranno prese in
considerazione.}
\renewcommand{\thefootnote}{\arabic{footnote}} 
\vskip.1cm
\centerline{\Large Quarta Prova di Autovalutazione}
\vskip1cm
\begin{enumerate}
\item Disegnare il grafico della funzione 
\[
f(x)=\frac{x+2}{x^2-1}
\]
Quante soluzioni ha l'equazione $f(x) = -1$ ? 
\item Trovare la distanza minima tra l'origine ed i punti sulla retta
passante per $(1,0)$ e $(0,1)$.  
\item Tra tutti i rettangoli con un vertice nell'origine ed un altro
sull' iperbole $y = \frac 3 x$, per $x > 0$, determinare quello di
perimetro minore.  
\item Si disegni l'insieme $A = \{(x,y) \in {\mathbb R}^2  \;|\;  xy
\geq 0\}$.  
\item Si calcoli $e$ con la precisione di due cifre decimali senza
l'ausilio di un calcolatore.
\item Si dimostri che
\[
1-e^{-x}\leq x
\]
per ogni $x\geq 0$.

\end{enumerate}
\end{document}
